%!TEX program = xelatex
\documentclass[a4paper,12pt]{article}
\usepackage{fontspec}\defaultfontfeatures{Ligatures=TeX}
\usepackage[a4paper,vmargin={3cm,3.5cm},hmargin={3cm,3cm}]{geometry}
%-----------------------------------------------------------------------------%
\usepackage{settings}
%-----------------------------------------------------------------------------%
\title{G6411 Recitation Notes: Time Series Basics}
\author{Jesse Chieh Chen\\\href{mailto:cc5229@columbia.edu}{cc5229@columbia.edu}}
\date{\today}
%-----------------------------------------------------------------------------%

\begin{document}

\maketitle

\noindent
Time series observations have a meaningful order.
We ask whether a single time series can reveal population moments and how dependence affects inference.
Our goal is to develop \gls{lln} and \gls{clt} analogous to those for \gls{iid} data.

\section{Terminology and Basic Properties}

\begin{definition}[Time Series]\label{def-time-series}
	A time series is a sequence of random variables $\{Y_t\}$ indexed by discrete time.
	We observe a sample $Y_1,\ldots,Y_T$ from this process.
\end{definition}

\begin{definition}[Time Series and Autocovariance]
	Let $\{Y_t\}$ be a time series with finite first and second moments.
	Its lag-$k$ autocovariance is
	\begin{align}
		\gamma_t(k)=\Cov(Y_t,Y_{t-k}).
	\end{align}
	When the variances are positive, the corresponding autocorrelation is
	\begin{align}
		\rho_t(k)=\frac{\gamma_t(k)}
		{\sqrt{\Var(Y_t)\Var(Y_{t-k})}}.
	\end{align}
\end{definition}

\begin{remark}
	The function $\gamma_t$ is an unconditional covariance:
	the subscript $t$ records the date, not conditioning on information available at that date.
\end{remark}

\begin{definition}[Covariance Stationarity]\label{def-cov-stationary}
	A time series is \emph{covariance stationary} (or \emph{weakly stationary}) if it has finite second moments and
	\begin{align}
		\E Y_t=\mu,\qquad \Cov(Y_t,Y_{t-k})=\gamma(k)
		\quad\text{for all valid dates and lags}.
	\end{align}
	When $\gamma(0)>0$, the \gls{acf} is $\rho(k)=\gamma(k)/\gamma(0)$.
\end{definition}

\begin{remark}
	If $\{Y_t\}$ is a scalar covariance-stationary process,
	then $\gamma(k)=\gamma(-k)$ and $\gamma(0)=\Var(Y_t)$.
	If $\{Y_t\}$ is a vector-valued covariance-stationary process,
	let $\Gamma(k)$ be the lag-$k$ autocovariance matrix.
	Then, $\Gamma(-k)=\Gamma(k)\T$ since $\Cov(Y_t,Y_{t-k})=\Cov(Y_{t-k},Y_t)\T$.
\end{remark}

\begin{example}[Random Walk]
	Let $e_t\iidto\normaldistribution(0,\sigma^2)$ with $\sigma^2>0$.
	For $Y_t=Y_{t-1}+e_t$ with fixed $Y_0=0$,
	\begin{align}
		Y_t=\sum_{j=1}^t e_j,\qquad \Var(Y_t)=t\sigma^2.
	\end{align}
	The mean is zero, but the variance changes with time.
\end{example}

\begin{definition}[Covariance Matrix of Covariance-Stationary Samples]
	A covariance-stationary sample $\{Y_t\}_{t=1}^{T}$ has covariance matrix
	\begin{align}
		\begin{pmatrix}
			\gamma(0) & \gamma(1) & \gamma(2) & \cdots & \gamma(T-1)\\
			\gamma(1) & \gamma(0) & \gamma(1) & \cdots & \gamma(T-2)\\
			\gamma(2) & \gamma(1) & \gamma(0) & \cdots & \gamma(T-3)\\
			\vdots    & \vdots    & \vdots    & \ddots & \vdots\\
			\gamma(T-1)&\gamma(T-2)&\gamma(T-3)&\cdots&\gamma(0)
		\end{pmatrix}.
	\end{align}
	When $\gamma(0)>0$, dividing this matrix by $\gamma(0)$ gives the correlation matrix.
\end{definition}

\begin{remark}[Toeplitz]
	A Toeplitz matrix is constant along each diagonal.
	The covariance matrix of a covariance-stationary process is Toeplitz.
	Some linear algebra results are easier to state for Toeplitz matrices than for general matrices.
\end{remark}

\begin{definition}[Strict Stationarity]\label{def-strict-stationary}
	A process is strictly stationary if, for every finite collection of dates $t_1,\ldots,t_m$
	and every integer shift $h$ for which all the dates are valid,
	\begin{align}
		(Y_{t_1},\ldots,Y_{t_m})\overset{d}{=}(Y_{t_1+h},\ldots,Y_{t_m+h}).
	\end{align}
\end{definition}

\begin{remark}[Strict versus Covariance Stationarity]
	\nameref{def-strict-stationary} with finite second moments implies \nameref{def-cov-stationary}.
	The converse fails.
	See \autoref{ex-same-moments}.
\end{remark}

\begin{example}[Identical Marginal Distributions Are Not Enough]
	Toss a fair coin independently each time and record each outcome twice.
	Writing $Z_j=1$ for heads and $Z_j=0$ for tails, the recorded sequence is
	\begin{align}
		Y_{2j-1}=Y_{2j}=Z_j,\qquad j=1,2,\ldots.
	\end{align}
	Every $Y_t$ equals zero or one with equal probability, so all marginal distributions are identical.
	However,
	\begin{align}
		\Pr(Y_1=Y_2)=1,\qquad \Pr(Y_2=Y_3)=\frac12.
	\end{align}
	The joint distribution of two consecutive observations changes when we shift their dates.
	Thus the process is not strictly stationary even though all marginal distributions are identical.
\end{example}

\begin{example}[Same Moments, Different Distributions]\label{ex-same-moments}
	An \gls{iid} standard normal process $\{e_t\}$ and an \gls{iid} process $\{u_t\}$ taking values $-1$ and $1$
	with equal probability both have mean zero and
	\begin{align}
		\gamma(0)=1,\qquad \gamma(k)=0\quad\forall k\neq0.
	\end{align}
	Assume the two processes are independent of each other,
	and define a new process $\{Y_t\}$ by the following:
	\begin{align}\label{eq-even-odd-process}
		Y_t=
		\begin{cases}
			e_t, & t\text{ even},\\
			u_t, & t\text{ odd}.
		\end{cases}
	\end{align}
	Then $\{Y_t\}$ is covariance stationary but not strictly stationary.
\end{example}

\begin{remark}[Second-Order Equivalence]
	There is this notion called \emph{second-order equivalence}:
	the processes have the same mean and autocovariance function.
	In \autoref{ex-same-moments}, $\{e_t\}$ and $\{u_t\}$ are second-order equivalent.
\end{remark}

\begin{definition}[White Noise]
	A process $\{e_t\}$ is \emph{white noise}, written $e_t\sim\mathrm{WN}(0,\sigma^2)$, if
	\begin{align}
		\E e_t=0,\qquad \Var(e_t)=\sigma^2<\infty,\qquad
		\Cov(e_t,e_s)=0\quad\text{ for }t\neq s.
	\end{align}
\end{definition}

\begin{remark}
	A white noise process is covariance stationary, but not necessarily strictly stationary.
	The white-noise assumption imposes only restrictions on means and covariances.
	For example, the process defined in \eqref{eq-even-odd-process} is white noise, but clearly not strictly stationary.
	A white noise process can even be dependent across time.
	See \autoref{ex-uncorrelated-dependent}.
\end{remark}

\begin{example}[Uncorrelated but Dependent]\label{ex-uncorrelated-dependent}
	Let $Z_t$ be \gls{iid} standard normal and set $e_t=Z_tZ_{t-1}$.
	Then $\E e_t=0$, $\Var(e_t)=1$, and all nonzero-lag autocovariances vanish.
	Therefore, $\{e_t\}$ is white noise.
	Nevertheless, adjacent observations share $Z_{t-1}$ and are dependent:
	\begin{align}
		\E(e_t^2e_{t-1}^2)
		=\E Z_t^2\,\E Z_{t-1}^4\,\E Z_{t-2}^2=3
		\neq \E e_t^2\,\E e_{t-1}^2=1.
	\end{align}
\end{example}

\begin{definition}[Correlogram]\label{def-correlogram}
	A correlogram is a graph where the $y$-axis is the autocorrelation $\rho(k)$ and the $x$-axis is the lag $k$.
\end{definition}

\begin{remark}
	\begin{figure}[t]
		\centering
		\includegraphics[width=\linewidth]{figures/acf.pdf}
		\caption{Examples of Correlograms.}
		\label{fig-acf}
	\end{figure}
	\autoref{fig-acf} illustrates three population \glspl{acf}:
	white noise has zero autocorrelation at every nonzero lag;
	gradual decay indicates that linear dependence persists across many lags;
	and a cutoff indicates zero autocorrelation beyond a certain lag.
\end{remark}

\section{Ergodicity and Ergodic Theorem}\label{sec-ergodicity}

\begin{remark}[What is Ergodicity?]
	Ergodicity is a weird concept to grasp.
	The idea is that it is the answer to the following question:
	\begin{quote}
		What condition allows \emph{time averages} to recover population averages without independence?
	\end{quote}
	We know that \gls{iid} is sufficient, but it is not necessary.
	\emph{Ergodicity} is a much weaker condition under which an \gls{lln} is guaranteed,
	and that \gls{lln} is typically known as the \emph{\nameref{thm-ergodic}}.
	In this course, we will refrain from defining ergodicity formally,
	as it is not worth our time to do so.
	Nevertheless, \autoref{ex-casino} provides an intuitive understanding of ergodicity.
\end{remark}

\begin{remark}
	An \gls{iid} process is ergodic.
\end{remark}

\begin{example}[Slot Machines]\label{ex-casino}
	Imagine a casino with $100$ identical slot machines.
	The machines operate independently,
	and the pattern of payouts follows the same distribution whenever we start watching.
	However, payouts depend on previous plays,
	since each machine starts spinning from where it stopped on the previous play.
	To decide whether playing is worthwhile,
	we want to know the \emph{expected payout per play}.
	There are two ways we might try to do this:
	\begin{enumerate}
		\item
			Play the $100$ machines once each and average their payouts.
			Since the machines operate independently, these payouts are \gls{iid}.
			The usual \gls{lln} says that this average converges to the expected payout as the number of machines grows.
			This is an \emph{ensemble} (or spatial) average.
		\item
			Playing each of the 100 machines once is very suspicious.
			How about playing \emph{one} machine repeatedly like a normal person and averaging its payouts?
			This is a \emph{time} average along one history generated by the machine's rule.
			As the number of plays grows, does it converge to the same expected payout
			that the average across machines estimates?
	\end{enumerate}
	If the payout sequence is \emph{ergodic},
	both the \emph{time average} and the \emph{ensemble average} are guaranteed to converge to the same \emph{expected payout per play}.
	Roughly speaking,
	an \emph{ergodic} sequence reflects the population distribution in its long-run frequencies,
	so one history is representative of the population.
\end{example}

\begin{definition}[Strong Mixing]
	For a process $\{Y_t\}_{t=1}^{\infty}$,
	let $\mathcal A_t$ be the collection of events
	determined by $\{Y_1,\ldots,Y_t\}$, and let $\mathcal B_{t+k}$ be the collection of events
	determined by $\{Y_{t+k},Y_{t+k+1},\ldots\}$.
	For $k\ge1$, define
	\begin{align}
		\alpha(k)&\coloneqq\sup_{t\ge1}\sup_{\substack{A\in\mathcal A_t\\B\in\mathcal B_{t+k}}}\alpha(A,B) \\
		\shortintertext{where}
		\alpha(A,B)&\coloneqq \left|\Pr(A\cap B)-\Pr(A)\Pr(B)\right|.
	\end{align}
	The function $\alpha(k)$ is known as the \emph{mixing coefficient} at lag $k$.
	The process is said to be \emph{strongly mixing} or \emph{$\alpha$-mixing} if $\alpha(k)\to0$ as $k\to\infty$.
\end{definition}

\begin{remark}
	For \gls{iid} observations, $\alpha(k)=0$ for every $k\ge1$.
\end{remark}

\begin{remark}[Interpreting Strong Mixing]
	Independence would give $\Pr(A\cap B)=\Pr(A)\Pr(B)$,
	so $\alpha(A,B)$ measures the
	``non-independent-ness'' between events $A\in\mathcal{A}_t$ and $B\in\mathcal{B}_{t+k}$.
	The coefficient $\alpha(k)$ measures the largest departure from this equality
	over all such events and all dates $t$.
	\emph{Strong mixing} says that any dependence between the past and the future vanishes uniformly as the time gap $k$ grows.
\end{remark}

\begin{proposition}
	A strictly stationary, strongly mixing process is ergodic.
\end{proposition}

\begin{remark}
	In a sense, ergodicity is also about ``mixing,'' but it is strictly weaker than strong mixing.
	In particular, ergodicity does not require dependence between observations to vanish as their temporal separation increases.
\end{remark}

\begin{example}[A Permanent Random Component]\label{ex-permanent-component}
	Let $Y_t=A+\varepsilon_t$, where $\E A=0$, $\Var(A)=\tau^2>0$,
	and $\varepsilon_t$ is \gls{iid}, independent of $A$, with mean zero and variance $\sigma^2$.
	This process is strictly stationary, but
	\begin{align}
		\overline Y_T=A+\frac1T\sum_{t=1}^T\varepsilon_t\pto A,\qquad \E Y_t=0.
	\end{align}
	Time averaging removes $\varepsilon_t$, not $A$.
	One history cannot average over the possible permanent components, so the process is not ergodic.
\end{example}

\begin{theorem}[Ergodic Theorem]\label{thm-ergodic}
	If $\{Y_t\}$ is strictly stationary and ergodic and $\E|Y_t|<\infty$, then
	\begin{align}
		\overline Y_T=\frac1T\sum_{t=1}^T Y_t\pto\mu=\E Y_t.
	\end{align}
	More generally, for fixed $m$ and an integrable function $h$,
	\begin{align}
		\frac1T\sum_{t=m+1}^T h(Y_t,\ldots,Y_{t-m})
		\pto \E h(Y_t,\ldots,Y_{t-m}).
	\end{align}
	The sum starts at $m+1$ so that all $m$ lags are observed.
\end{theorem}

\begin{corollary}[Estimating Autocovariances]\label{cor-sample-acf}
	If the process is strictly stationary and ergodic with $\E Y_t^2<\infty$, then for every fixed $k\ge0$,
	\begin{align}
		\widehat\gamma_T(k)
		=\frac1T\sum_{t=k+1}^T(Y_t-\overline Y_T)(Y_{t-k}-\overline Y_T)\pto\gamma(k).
	\end{align}
	If $\gamma(0)>0$, then $\widehat\rho_T(k)=\widehat\gamma_T(k)/\widehat\gamma_T(0)\pto\rho(k)$.
\end{corollary}
\begin{proof}
	Apply \nameref{thm-ergodic} to $Y_t$ and to $Y_tY_{t-k}$.
	The latter is integrable by Cauchy-Schwarz.
	Expanding the centered product and using $\overline Y_T\pto\mu$ gives the result.
\end{proof}

\section{\glsentrylong{clt} for Dependent Data}\label{sec-clt-dependent}

Let $\{Y_t\}$ be a covariance-stationary process with mean $\mu$ and autocovariance function $\gamma(k)$.
Consider the sample mean $\overline Y_T=T\inv\sum_{t=1}^T Y_t$:
\begin{align}
	\Var(\overline Y_T)
	=\frac1{T^2}\sum_{t=1}^T\sum_{s=1}^T\gamma(t-s)
	&=\frac1T\sum_{k=-(T-1)}^{T-1}\left(1-\frac{|k|}{T}\right)\gamma(k)\nonumber\\
	&=\frac1T\left[\gamma(0)+2\sum_{k=1}^{T-1}
	\left(1-\frac{k}{T}\right)\gamma(k)\right].\label{eq-mean-var}
\end{align}

\begin{remark}
	Positive autocovariances increase the variance of the mean, whereas negative ones reduce it.
	Hence, $T$ dependent observations need not be more precise or imprecise than $T$ \gls{iid} observations.
	If $\{Y_t\}$ is \gls{iid}, then $\Var(\overline Y_T)=\gamma(0)/T$.
\end{remark}

\begin{definition}[Absolutely Summable Covariance]\label{def-absolutely-summable}
	A covariance-stationary process $\{Y_t\}$ has \emph{absolutely summable covariance} if
	\begin{align}
		\sum_{k=-\infty}^{\infty}|\gamma(k)|<\infty.
	\end{align}
	For a vector-valued process $Y_t\in\reals^k$,
	its autocovariances are absolutely summable if the above holds for every entry of the autocovariance matrix $\Gamma(k)$.
\end{definition}

\begin{definition}[\glsentrylong{lrv}]\label{def-lrv}
	Let $\{Y_t\}$ be a covariance-stationary process in $\reals^k$ with \nameref{def-absolutely-summable}.
	Then
	\begin{align}
		\Omega_T\coloneqq T\Var(\overline Y_T)
		&=\Gamma(0)+\sum_{j=1}^{T-1}\left(1-\frac jT\right)
		[\Gamma(j)+\Gamma(j)\T]\tag{\ref{eq-mean-var}$'$}\\
		&\to\Omega\coloneqq\sum_{j=-\infty}^{\infty}\Gamma(j).\label{eq-lrv}
	\end{align}
	The limit $\Omega$ is called the \gls{lrv} of the process.
\end{definition}
\begin{proof}[Proof of Convergence]
	Expanding covariances as in \eqref{eq-mean-var} gives
	\begin{align*}
		\Omega_T=\sum_{|j|<T}\left(1-\frac{|j|}{T}\right)\Gamma(j).
	\end{align*}
	Since $\Gamma(-j)=\Gamma(j)\T$, this yields the displayed formula.
	To prove convergence, fix an entry $(a,b)$ and any $\varepsilon>0$.
	Absolute summability lets us choose $K$ such that
	$\sum_{|j|>K}|\Gamma_{ab}(j)|<\varepsilon/2$.
	For $T>K$,
	\begin{align*}
		|(\Omega_T)_{ab}-\Omega_{ab}|
		\le\sum_{|j|<T}\frac{|j|}{T}|\Gamma_{ab}(j)|
		+\sum_{|j|\ge T}|\Gamma_{ab}(j)|
		\le\frac KT\sum_{|j|\le K}|\Gamma_{ab}(j)|
		+\sum_{|j|>K}|\Gamma_{ab}(j)|.
	\end{align*}
	With this $K$ fixed, choose $T$ large enough that the first term is also less than $\varepsilon/2$.
	Thus every entry converges, so $\Omega_T\to\Omega$.
\end{proof}

\begin{remark}[Some Examples]\leavevmode
	\begin{itemize}
		\item
			For a scalar process, we have $\omega_T=T\Var(\overline Y_T)$ and
			$\omega=\sum_{j=-\infty}^{\infty}\gamma(j)$.
		\item
			For white noise or \gls{iid} data, we have $\omega_T=\omega=\gamma(0)=\sigma^2$.
		\item
			For the process defined in \autoref{ex-permanent-component},
			$\Var(\overline Y_T)=\tau^2+\sigma^2/T$ and $\omega_T=T\tau^2+\sigma^2$ diverges as $T\to\infty$.
			Here, $\gamma(j)=\tau^2>0$ for every nonzero lag.
	\end{itemize}
\end{remark}

\begin{theorem}[\gls{clt} for Dependent Data]\label{thm-clt-correlated}
	Let $\{Y_t\}$ be a strictly stationary process in $\reals^k$ with mean $\mu$
	and mixing coefficients $\alpha(k)$.
	Suppose that for some $r>2$, $\E\|Y_t\|^r<\infty$ and $\sum_{j=1}^{\infty}\alpha(j)^{1-2/r}<\infty$.
	Then
	\begin{align}
		\sqrt T(\overline Y_T-\mu)\dto\normaldistribution(0,\Omega),
	\end{align}
	where $\Omega$ is the \gls{lrv} of $\{Y_t\}$.
\end{theorem}

\begin{remark}[Ergodicity is Not Enough]
	The mixing condition in \autoref{thm-clt-correlated} is much stronger than $\alpha(k)\to0$ as $k\to\infty$.
	Ergodicity alone guarantees the \gls{lln} provided that the mean exists,
	but does not generally guarantee a \gls{clt}.
\end{remark}

\begin{remark}
	\autoref{thm-clt-correlated},
	which is Theorem 14.15 in \textcite{hansen-econometrics},
	is but one of many \glspl{clt} for dependent data,
	just like Lindeberg-Lévy \gls{clt} is but one of many \glspl{clt} for independent data.
\end{remark}

\begin{remark}
	Once a \gls{clt} is established,
	the usual normal approximations and delta-method arguments apply.
	To make these feasible, we need a consistent estimator of the \gls{lrv} $\Omega$.
\end{remark}

\subsection{Estimation of the \glsentrylong{lrv}}

\begin{remark}[Why Not Estimate Every $\gamma(k)$?]
	For a scalar process, under suitable moment and dependence conditions,
	each fixed-lag estimator has variance $O(1/T)$ and estimation error $O_p(T^{-1/2})$.
	But the number of available lags also grows with $T$.
	Even if the estimation errors were uncorrelated,
	summing order $T$ terms with variances of order $1/T$ would leave a variance of order one.
	Thus, consistency at each fixed lag alone does not guarantee consistency of the growing sum.
\end{remark}

\begin{definition}
	For $Y_t\in\reals^k$, define the sample autocovariance matrix as
	\begin{align}
		\widehat\Gamma_T(j)=\frac1T\sum_{t=j+1}^T (Y_t-\overline Y_T)(Y_{t-j}-\overline Y_T)\T,
		\qquad 0\le j<T.
	\end{align}
	For $1\le j<T$, define $\widehat\Gamma_T(-j)=\widehat\Gamma_T(j)\T$.
	This is the vector version of the estimator in \autoref{cor-sample-acf}.
\end{definition}

\begin{remark}[Truncating the Sum]
	Choose a maximum lag $L<T$ and consider
	\begin{align}\label{eq-lrv-truncated}
		\widehat\Omega_T^{\mathrm{cut}}(L)
		=\widehat\Gamma_T(0)+\sum_{j=1}^{L}
		[\widehat\Gamma_T(j)+\widehat\Gamma_T(j)\T].
	\end{align}
	For fixed $L$, this converges to $\sum_{j=-L}^{L}\Gamma(j)$, not generally to $\Omega$.
	It works if the omitted covariances sum to zero, for example if all autocovariances vanish beyond $L$.
	Otherwise, we must let $L=L_T\to\infty$, but slowly enough that estimation error does not accumulate too quickly.
\end{remark}

\begin{remark}[Weighting the Lags]
	A growing cutoff can produce a consistent estimator under suitable conditions,
	but the abrupt cutoff has a drawback: the estimate need not be positive semidefinite.
	See \autoref{ex-negative-variance}.
\end{remark}

\begin{example}\label{ex-negative-variance}
	Take the scalar sample $(Y_1,Y_2,Y_3)=(1,-2,1)$, so $\overline Y_3=0$.
	Then
	\begin{align*}
		\widehat\gamma_3(0)=\frac{1+4+1}{3}=2,
		\qquad \widehat\gamma_3(1)=\frac{-2-2}{3}=-\frac43.
	\end{align*}
	With cutoff $L=1$, the \gls{lrv} estimate is
	$\widehat\gamma_3(0)+2\widehat\gamma_3(1)=-2/3$.
\end{example}

\begin{definition}[Kernel]
	A \emph{kernel} $K(j,L)$ assigns a weight to lag $j$ given a nonnegative integer \emph{bandwidth} $L$.
	We consider kernels satisfying
	\begin{enumerate}
		\item(\textsc{Symmetry}) $K(-j,L)=K(j,L)$ for all $j$ and $L$.
		\item(\textsc{Boundedness}) $\exists C<\infty$ such that $|K(j,L)|\le C$ for all $j$ and $L$.
		\item(\textsc{Compact Support}) $K(j,L)=0$ for $|j|>L$.
		\item(\textsc{Normalization}) $K(0,L)=1$ for all $L$.
		\item(\textsc{Fixed-Lag Convergence}) $K(j,L)\to1$ as $L\to\infty$ for every fixed $j$.
	\end{enumerate}
\end{definition}

\begin{remark}[Properties of Kernels]
	A kernel is simply a rule for choosing weights on the sample autocovariances.
	These conditions keep the weights bounded and symmetric, leave lag zero unchanged,
	and omit lags beyond the bandwidth.
	As the bandwidth grows, each fixed lag eventually receives weight close to one,
	so its contribution is not permanently discounted.
\end{remark}

\begin{definition}[Kernel Estimator of the \gls{lrv}]
	Given a kernel $K$ and a bandwidth $L_T<T$, define
	\begin{align}\label{eq-hac}
		\widehat\Omega_T
		= \sum_{j=-(T-1)}^{T-1} K(j,L_T)\widehat\Gamma_T(j)
		= \widehat\Gamma_T(0)+\sum_{j=1}^{L_T} K(j,L_T) [\widehat\Gamma_T(j)+\widehat\Gamma_T(j)\T].
	\end{align}
\end{definition}

\begin{remark}
	The uniform kernel $K(j,L)=\mathbf{1}\{|j|\le L\}$ gives the cutoff in \eqref{eq-lrv-truncated}.
\end{remark}

\begingroup
\tcbset{lines before break=4}
\begin{definition}[Bartlett Kernel and \gls{nw} Estimator]
	The \emph{Bartlett kernel} uses triangular weights:
	\begin{align}
		K(j,L)=\left(1-\frac{|j|}{L+1}\right)\mathbf{1}\{|j|\le L\}.
	\end{align}
	Substituting it into \eqref{eq-hac} gives the \emph{\gls{nw} estimator}
	\begin{align}\label{eq-newey-west}
		\widehat\Omega_T^{\mathrm{NW}}
		=\widehat\Gamma_T(0)+\sum_{j=1}^{L_T}\left(1-\frac{j}{L_T+1}\right)
		[\widehat\Gamma_T(j)+\widehat\Gamma_T(j)\T].
	\end{align}
\end{definition}
\endgroup

\begin{remark}[Positive Semidefiniteness]
	A useful property of Bartlett weights is that the \gls{nw} estimator is positive semidefinite for \textsc{every sample}.
	In particular, its scalar version cannot be negative, unlike the abrupt cutoff estimator.
\end{remark}

\begin{example}
	Consider again \autoref{ex-negative-variance} with $L=1$.
	Then the \gls{nw} estimate is
	\begin{align}
		\widehat\gamma_3(0)+\frac12[\widehat\gamma_3(1)+\widehat\gamma_3(1)]
		=2+\frac12\left(-\frac43-\frac43\right)=\frac23>0.
	\end{align}
\end{example}

\begin{theorem}[Consistency of the \gls{nw} Estimator]\label{thm-nw-consistency}
	Let $\{Y_t\}$ be strictly stationary and ergodic with mixing coefficients $\alpha(j)$.
	Suppose that, for some $r>4$,
	$\E\|Y_t\|^r<\infty$ and
	$\sum_{j=1}^{\infty}\alpha(j)^{1/2-2/r}<\infty$.
	If $L_T\to\infty$ and $L_T^3/T=O(1)$, then
	\begin{align}
		\widehat\Omega_T^{\mathrm{NW}}\pto\Omega.
	\end{align}
\end{theorem}

\begin{remark}
	\autoref{thm-nw-consistency} is adapted from Theorem 14.34 in \textcite{hansen-econometrics}.
	Hansen imposes moment conditions on the regression variables, whereas his \gls{nw} estimator uses the products $X_te_t$.
	By Cauchy-Schwarz, finite moments of order $2r$ for $X_t$ and $e_t$ imply finite moments of order $r$ for $X_te_t$.
	Thus, the corresponding moment order in Hansen's notation is $2r$, and his mixing exponent becomes $1/2-4/(2r)=1/2-2/r$.
\end{remark}

\begin{remark}[Choosing the Bandwidth: \glsentryshort{hac}/\glsentryshort{har}]
	Traditional \gls{hac} inference uses a consistent estimator of the \gls{lrv},
	with $L_T\to\infty$ but $L_T/T\to0$, subject to additional rate conditions.
	For example, $L_T=\lfloor cT^{1/3}\rfloor$ with fixed $c>0$ satisfies the preceding theorem.
	Consistency justifies the usual normal or chi-squared critical values asymptotically,
	but finite-sample rejection rates can still be far from the nominal size.

	\gls{har} inference focuses on the validity of tests rather than requiring a consistent variance estimator.
	One approach is \emph{fixed-$b$} inference, which takes $L_T=\lfloor bT\rfloor$ for fixed $0<b<1$.
	The variance estimator is then generally inconsistent, but under suitable conditions
	tests remain asymptotically valid using nonstandard critical values.
\end{remark}

\begin{example}[Linear Regression]\label{ex-linear-regression}
	Let $\{(Y_t,X_t)\}$, with $Y_t\in\reals$ and $X_t\in\reals^k$, be a strictly stationary time series satisfying
	\begin{align*}
		Y_t=X_t\T\beta+e_t
		\quad\text{where}\quad
		\E(X_te_t)=0.
	\end{align*}
	Let $Q_{xx}\coloneqq \E(X_tX_t\T)$ be finite and positive definite.
	Let $g_t\coloneqq X_te_t$ have absolutely summable covariances and $\Omega_g\coloneqq \sum_{j=-\infty}^{\infty}\Gamma_g(j)$ be the \gls{lrv} of $\{g_t\}$.
	If $\{X_t\}$ is ergodic, then
	\begin{align}
		\widehat Q_{xx}\coloneqq\frac1T\sum_{t=1}^T X_tX_t\T\pto Q_{xx}.
	\end{align}
	Furthermore, if the conditions of \autoref{thm-clt-correlated} hold for $\{g_t\}$,
	then the \gls{ols} estimation error satisfies
	\begin{align*}
		\sqrt T(\widehat\beta-\beta)
		=\widehat Q_{xx}\inv\frac1{\sqrt T}\sum_{t=1}^T X_te_t
		\dto\normaldistribution(0,V)
		\quad\text{where}\quad
		V=Q_{xx}\inv\Omega_gQ_{xx}\inv.
	\end{align*}
	To estimate $V$, we need a consistent estimator of $\Omega_g$.
	Estimate the score autocovariances by
	\begin{align*}
		\widehat\Gamma_{g,T}(j)
		\coloneqq\frac1T\sum_{t=j+1}^T\widehat g_t\widehat g_{t-j}\T,
		\qquad 0\le j\le L_T,
	\end{align*}
	where $\widehat g_t\coloneqq X_t\widehat e_t=X_t(Y_t-X_t\T\widehat\beta)$.
	The \gls{ols} normal equations give $\sum_t\widehat g_t=0$, so no further centering is needed.
	Now we can apply \gls{nw} according to \eqref{eq-newey-west} to estimate the \gls{lrv} of the scores.
	If the conditions of \autoref{thm-nw-consistency} hold for $\{g_t\}$ and $\E\|X_t\|^4<\infty$,
	then the feasible \gls{nw} estimator is consistent:
	\begin{align*}
		\widehat\Omega_{g,T}^{\mathrm{NW}}
		=\widehat\Gamma_{g,T}(0)+\sum_{j=1}^{L_T}\left(1-\frac{j}{L_T+1}\right)
		[\widehat\Gamma_{g,T}(j)+\widehat\Gamma_{g,T}(j)\T]
		\pto
		\Omega_g.
	\end{align*}
	Thus $\widehat V_{\mathrm{NW}}=\widehat Q_{xx}\inv\widehat\Omega_{g,T}^{\mathrm{NW}}\widehat Q_{xx}\inv\pto V$.
	If $L_T=0$, this sandwich reduces to HC0.
	This type of procedure applies to many time series models that we will see later.
\end{example}

\begin{remark}[Using Estimated Scores]\label{rem-estimated-scores}
	As with HC0, consistency with the true scores $g_t$ does not automatically justify using $\widehat g_t$.
	We must also control estimation error over the included lags;
	the additional condition $\E\|X_t\|^4<\infty$ suffices in \autoref{ex-linear-regression}.
\end{remark}

\vfill
\section*{Acronyms}
\printglossaries

\section*{References}
\printbibliography[heading=none]

\end{document}
